% The next line makes it so it compiles the main file if I hit the compile button
% when editing this file in TeXworks.
% !TEX root = ../ActiveIntroToDiscreteMathAndAlgorithms.tex
\section{Problems}

\begin{pro}\label{pro:booExp}
Draw a truth table to represent the following.
\begin{enumerate}[(a)]
\item $\neg p \vee q$
\item $(p\rightarrow q) \vee \neg p$
\item $(p\wedge \neg q)\vee r$
\item $((p\vee q) \wedge\neg(p\vee q))\vee r$
\item $(p\vee\neg r)\wedge q$
\item $(p\oplus q)\wedge (q\vee r)$
\item $p\leftrightarrow (p\wedge q)$
\end{enumerate}
\end{pro}

\begin{pro}
Prove the distributive laws.
\begin{enumerate}[(a)]
\item
$p\wedge(q\vee r) = (p\wedge q) \vee (p\wedge r)$
\item
$p\vee(q\wedge r) = (p\vee q)\wedge(p\vee r)$
\end{enumerate}
\end{pro}

\begin{pro}
Prove the identity laws.
\begin{enumerate}[(a)]
\item
$p \vee F = p$ 
\item
 $p\wedge T = p$
\end{enumerate}
\end{pro}

\begin{pro}
Prove the negation laws.
\begin{enumerate}[(a)]
\item
$p\vee \neg p = T$ 
\item
 $p\wedge\neg p = F$
\end{enumerate}
\end{pro}

\begin{pro}
Prove that $p\oplus q = (p\vee q)\wedge \neg(p\wedge q)$.
\end{pro}

\begin{pro}
Prove the following laws involving implications.
\begin{enumerate}[(a)]
\item
$p\rightarrow q = \neg p \vee q$
\item
$p\rightarrow q = \neg q \rightarrow \neg p$
\end{enumerate}
\end{pro}

\begin{pro}
Prove the following laws involving biconditionals.
\begin{enumerate}[(a)]
\item
$p\leftrightarrow q = (p\rightarrow q)\wedge (q\rightarrow p)$
\item
$p\leftrightarrow q = \neg p\leftrightarrow \neg q$
\item
$\neg(p\leftrightarrow q) = p\leftrightarrow\neg q$
\end{enumerate}
\end{pro}

\begin{pro}
Give 2 different proofs that $[(p\vee q)\wedge \neg p] \rightarrow q$ 
is a tautology. 
\end{pro}

\begin{pro}
Prove $\neg(p\leftrightarrow q) = p\oplus q$ without using truth tables.
\end{pro}

\begin{pro}
Express $p\vee q \vee r$ using only $\wedge$ and $\neg$ .
\end{pro}

\begin{pro}
 The {\em NAND} of $p$ and $q$, denoted by $p | q$, is the proposition ``not both $p$ and $q$''. 
The NAND of $p$ and $q$ is false when $p$ and $q$ are both true and true otherwise.
\begin{enumerate}[(a)]
\item Draw a truth table for $NAND$
\item Express $p|q$ using $\vee$, $\wedge$, and/or $\neg$ (you may not need all of them).
\item Express $p\wedge q$ using {\em only} $|$. 
(That means you cannot use $\neg$, $\vee$, $\wedge$, or any other
boolean operator except for $|$. Thus, your answer should {\em only} involve $p$, $q$, $|$ and parentheses.)
Your answer should be as simple as possible. Give a truth table that shows they are the same.
\item Express $\neg p\vee q$ using only $|$.  Your answer should be as simple as possible. Give a truth table that shows they are the same.
\end{enumerate}
\end{pro}

\begin{pro}
 The {\em NOR} of $p$ and $q$, denoted by $p \downarrow q$, is the proposition ``neither  $p$ nor $q$''. 
The NOR of $p$ and $q$ is true when $p$ and $q$ are both false and false otherwise.
\begin{enumerate}[(a)]
\item Draw a truth table for $\downarrow$
\item Express $p\downarrow q$ using $\vee$, $\wedge$, and/or $\neg$ (you may not need all of them).
\item Express $p\wedge q$ using {\em only} $\downarrow$. 
(That means you cannot use $\neg$, $\vee$, $\wedge$, or any other
boolean operator except for $\downarrow$. Thus, your answer should {\em only} involve $p$, $q$, $\downarrow$ and parentheses.)
Your answer should be as simple as possible.  Give a truth table that shows they are the same.
\item Express $\neg p \vee q$ using only $\downarrow$. Your answer should be as simple as possible.  Give a truth table that shows they are the same.
\end{enumerate}
\end{pro}

\begin{pro}
A set of logical operators is {\em functionally complete} if any possible operator can be implemented using only 
operators from that set.  It turns out that $\{\neg,\wedge\}$ is functionally complete.  So is  $\{\neg,\vee\}$.
To show that a set if functionally complete, all one needs to do is show how to implement all of the operators
from another functionally complete set.  Given this,
\begin{enumerate}[(a)]
\item
Show that $\{|\}$ is functionally complete.
(Hint: Since $\{\neg,\wedge\}$ is functionally complete, 
one way is to show how to implement both $\wedge$ and $\neg$ using just $|$.)
\item
Show that $\{\downarrow \}$ is functionally complete. 
\end{enumerate}
\end{pro}


\begin{pro}
Express the following phrase using quantifiers.
``There is some constant $c$ such that 
$f(x)$ is no greater than $c\cdot g(x)$ for all $x\geq x_0$ for some constant $x_0$.''
Your solution should contain no English words.
\end{pro}

\begin{pro}
Write each of the following expressions so that negations are only applied to
propositional functions (and not quantifiers or connectives).
\begin{lenum}
\item
$\neg \forall x \exists y \neg P(x,y)$
\item
$\neg (\forall x \exists y P(x,y) \wedge \exists x \neg \forall y P(x,y) )$
\item
$\neg \forall x (\exists y P(x,y) \vee \forall y Q(x,y))$
\item
$\neg \forall x \neg \exists y (\neg \forall z P(x,z) 
\rightarrow \exists z Q(x,y,z))$
\item
$\neg \exists x (\neg \forall y [\exists z ( P(y,x,z) \wedge P(y,z,x)\wedge 
P(x,y,z) ) ] \vee \exists z Q(x,z) )$
\end{lenum}
\end{pro}

\begin{pro}
Let $P(x,y)$=``$x$ likes $y$'', where the universe of discourse for
$x$ and $y$ is the set of all people.
Translate each of the following into English, smoothing them out as much as possible.
Then give the truth value of each.
\begin{lenum}
\item
$\forall x\forall y P(x,y)$
\item
$\forall x\exists y P(x,y)$
\item
$\forall y\exists x P(x,y)$
\item
$\forall x P(x, Raymond)$
\item
$\neg \forall x \forall y P(x,y)$
\item
$\forall x \neg \forall y P(x,y)$
\item
$\forall x \neg \forall y \neg P(x,y)$
\end{lenum}
\end{pro}

\begin{pro}
Let $P(x,y,z)$=``$x^2+y^2=z^2$'', where the universe of discourse for
all variables is the set of integers. 
What are the truth values of each of the following?
\begin{lenum}
\item
$\forall x\forall y\forall z P(x,y,z)$
\item
$\exists x\exists y\forall z P(x,y,z)$
\item
$\forall x\exists y\exists z P(x,y,z)$
\item
$\forall x\forall y\exists z P(x,y,z)$
\item
$\forall x\exists y\forall z P(x,y,z)$
\item
$\exists x\exists y\exists z P(x,y,z)$
\item
$\exists z P(2,3,z)$
\item
$\exists x \exists y P(x,y,5)$
\item
$\exists x \exists y P(x,y,3)$
\end{lenum}
\end{pro}

\begin{pro}
Write each of the following sentences using quantifiers and 
propositional functions.  Define propositional functions as necessary 
(e.g. Let $D(x)$ be the proposition `$x$ plays disc golf.')
\begin{lenum}
\item
All disc golfers play ultimate Frisbee.
\item
If all students in my class do their homework, then some of the
students will pass.
\item
If none of the students in my class study, then all of the students in
my class will fail.
\item
Not everybody knows how to throw a Frisbee 300 feet.
\item
Some people like ice cream, and some people like cake, but everybody
needs to drink water.
\item
Everybody loves somebody.
\item
Everybody is loved by somebody.
\item
Not everybody is loved by everybody.
\item
Nobody is loved by everybody.
\item
You can't please all of the people all of the time, but you can please
some of the people some of the time.
\item
If only somebody would give me some money, I would buy a new house.
\item
Nobody loves me, everybody hates me, I'm going to eat some worms.
\item
Every rose has its thorn, and every night has its dawn.
\item
No one ever is to blame.
\end{lenum}
\end{pro}

\begin{pro}
Consider the following expression: 
$$\forall\epsilon{>}0\, \exists \delta{>}0\, \forall x(0<|x-c|<\delta \rightarrow |f(x)-L|<\epsilon).$$
\begin{enumerate}[(a)]
\item
Express it in English.  Be as concise as possible.
\item
(Difficult if you have not had calculus.) This is the definition of something.  What is it?
\end{enumerate}
\end{pro}


\begin{pro}
Use Procedure~\ref{dnfConvert} to find the disjunctive normal form for each of the expressions from 
Problem \ref{pro:booExp}.
\end{pro}


